\documentclass[a4paper, 12pt]{article}

\usepackage[frenchb]{babel} % Document en français
\usepackage[utf8]{inputenc} % Document au format UTF8
\usepackage[T1]{fontenc} % Suppression d'un warning pour la langue francaise
\usepackage[left=2.2cm, right=2.2cm, top=2.5cm, bottom=2.5cm]{geometry} % Mise en page
\usepackage{setspace} % Permet de définir l'interligne
\usepackage{titling}
\usepackage{graphicx}
\usepackage{todonotes}

\singlespacing % Interligne à 1

\newcommand{\HRule}{\rule{\linewidth}{0.5mm}}
\graphicspath{{res/}}

\usepackage{newtxtext,newtxmath} % pour la police time


\begin{document}

\begin{titlepage}
\begin{center}

% Upper part of the page. The '~' is needed because \\
% only works if a paragraph has started.
\includegraphics[width=0.15\textwidth]{./sedes_ucl.png}~\\[1cm]

\textsc{\LARGE Université Catholique de Louvain}\\[1.5cm]

\textsc{\Large Travail pratique de Management humain}\\[0.5cm]

% Title
\HRule \\[0.4cm]
{ \huge \bfseries Les rôles de la fonction Ressources Humaines\\[0.4cm] }

\HRule \\[1.5cm]

% Author and supervisor
\noindent
\begin{minipage}[t]{0.4\textwidth}
\begin{flushleft} \large
\emph{Auteurs:}\\
Ivan \textsc{Ahad} \\
Aymeric \textsc{De Cocq} \\
Solaiman \textsc{El Jilali} \\
Léonard \textsc{Julémont} \\
Eddy \textsc{Ndizera} \\
Florian \textsc{Thuin}

\end{flushleft}
\end{minipage}%
\begin{minipage}[t]{0.4\textwidth}
\begin{flushright} \large
\emph{Professeurs :} \\
Nathalie \textsc{Delobbe} \\
Emilie \textsc{Malcourant}
\end{flushright}
\end{minipage}

\vfill

% Bottom of the page
{\large \today}

\end{center}
\end{titlepage}

%\maketitle

%\tableofcontents

\section{Présentation des différents conférenciers}
Notre groupe a eu l'occasion de participer à la table ronde accueillant trois directeurs de ressources humaines : Christine Thiran, Frédéric Nolf, Alain Goergen.

La première est Christine Thiran qui travaille actuellement comme directrice des ressources humaines pour le groupe Mestdagh. Durant son parcours, elle a occupé le poste de DRH chez Belgacom \todo{Dire le nombre exact d'années ou supprimer} durant X années  où elle a notamment été responsable d'un plan de restructuration concernant 4500 personnes \cite{planBEST}.
Elle a ensuite travaillé pendant 11 ans aux cliniques universitaires de Saint-Luc où elle a mis en place une section ressources humaines.
Depuis cinq mois, elle travaille au sein de cette entreprise familiale du secteur de la grande distribution. Ce secteur en grande difficulté disposant de marges limitées emploie majoritairement du personnel non-qualifié. Le groupe Mestdagh a comme vision d'offrir de l'emploi dans des zones économiquement défavorisées, et doit actuellement intégrer 16 anciennes enseignes (et leurs employés) du groupe Carrefour.
Durant cette courte période, Mme Thiran a émis prématurement un rapport d'étonnement questionnant entre autres les habitudes pour accélérer la prise de conscience des éléments à modifier.

Ce faisant, Mme Thiran s'est identifiée à plusieurs rôles dont la gestion d'intégration des employés des anciennes enseignes Carrefour rachetées. Elle remarqua un clivage persistant entre les employés Carrefour et Mestdagh. Elle doit donc avoir un rôle de \textbf{fédérateur} et rassembler l'ensemble des employés Mestdagh autour d'une vision commune, d'une culture commune. \todo{Pe a completer} garder l'esprit familial de l'entreprise et sa proximité avec les gens, offrir des perspectives de carrière au sein des différents magasins\newline  recrutement?


Parallèlement, Mme Thiran se considère comme un \textbf{moteur de changement}. Mestdagh évolue dans un secteur très compétitif et où les marges de bénéfices sont très petites. De ce fait, cela a un impact sur l'organisation et donc sur la gestion des ressources humaines. De ce fait, son rôle est de secouer l'ensemble de ses collaborateurs. Via ce rapport, elle tente d'initier un changement, une évolution de la situation chez Mestdagh. \newline
Pour ce faire, une de ses pratiques est de pousser l'ensemble des employés mais aussi des syndicats à participer, à etre consulté dans les prises de décisions. De plus, elle se voit comme le bras droit du CEO et veut entrer en symbiose avec celui-ci dans les prises de décisions. (Position optimale car proche du CEO et des employés + syndicats  ?)

Le second était Frédéric Nolf, responsable des ressources humaines pour IBA depuis 2007. IBA est une société côtée en bourse, basée sur l'innovation. Lors de sa création c'était une spin-off de l'UCL et elle est actuellement encore en alliance avec les universités pour trouver des nouveaux talents. Elle est côtée en bourse. Son CEO a changé il y a trois ans.

Son rôle est d'assurer à l'entreprise qu'elle dispose des bonnes personnes pour son développement et que les personnes soient dans la bonne entreprise pour leur développement. Ce rôle s'articule autour de plusieurs missions : s'assurer de la motivation de son personnel, assurer un \og{} sens du focus\fg{}(nous y reviendrons), créer des liens de compétence à travers les divers départements d'IBA.


Dans le cadre de ces missions, il doit relever de nombreux défis : faire rêver les ingénieurs nouvellement diplomés de travailler chez IBA, s'assurer que ces ingénieurs s'intègrent à la culture IBA, s'assurer que les ingénieurs travaillent réellement sur ce qui leur est demandé (\og{} quand il y a beaucoup de gens très intelligents, chacun a son idée sur ce qu'il faut faire \fg{}), informer les membres du personnel à travers le monde sur les projets actuels pour trouver les membres les plus aptes à les résoudre.

Ces défis nécessitent de mettre en place certaines pratiques de GRH : faire des enquêtes de satisfaction et y répondre de manière concrète, mettre en place des activités communes (sport), prendre en compte toutes les idées proposées, créer des projets qui lient les convictions personnelles du personnel à la culture de la société (mise en place d'une cellule de réflexion pour le développement durable).

\documentclass[12pt]{article}
%Pour les marges
\usepackage[a4paper]{geometry}
\geometry{hscale=0.75,vscale=0.86,centering} 
\usepackage[utf8]{inputenc}
\usepackage[frenchb]{babel}
\usepackage[T1]{fontenc}
\usepackage{amsmath}
\usepackage{amsfonts}
\usepackage{amssymb} 
\usepackage{graphicx}
\usepackage{lmodern}
\usepackage{numprint}
\usepackage{hyperref}
\usepackage{tabularx}
\usepackage[squaren, Gray]{SIunits}
\usepackage{ulem} %Pour barrer
\usepackage{listings}
\usepackage{enumitem}


\usepackage{newtxtext,newtxmath} % pour la police time

\title{Travail Management}
\author{Aymeric De Cocq, Léonard Julémont}
\date{\today}

\begin{document}

\section{Remise en contexte}


Alain Goergen est le directeur des ressources humaines de la Police Fédérale. Sa carrière s'est construite au sein de la police. Il commence par suivre une formation d'officier de gendarmerie à l'école militaire avec une licence en criminologie. Après 2 ans sur le terrain il entame une seconde formation en management pour se réorienter vers un poste administratif. Il a été en charge de nombreux projets toujours au sein des ressources humaines comme le recrutement, la formation ou encore la restructuration complète du département RH.\\

La Police Fédérale est le plus grand employeur de Belgique. Elle compte 193 employeurs qui emploient plus 50 000 personnes réparties dans différents domaines d'activités sur tout le territoire national. L'éventail de missions de la Police est assez large: maintien de l'ordre, assistance aux victimes, missions anti-terroristes, etc.\\
Malgré la division entre les 192 polices locales et la police fédérale, elles sont toutes dépendantes du département RH de la police fédérale.\\
 
La Police est aujourd'hui dans un phase difficile du point de vue économique et politique. En effet, le gouvernement doit faire des économies et ce service public n'est pas épargné. Le budget de la Police est donc moindre que les années précédentes [\textbf{trouver chiffres?}], pour des résultats identiques, voir même supérieurs. Nous pouvons illustrer ceci par les paroles d'Alain Goergen : 

\begin{center}
	\textit{Il y a de grosse difficulté budgétaire [...] ce qui implique une optimalisation. Faire aussi bien voir mieux mais avec beaucoup moins de moyens.}
\end{center} 

De plus, du fait des récents évènements terroristes en France, l'armée Belge travaille en coopération avec la police pour assurer la sécurité des citoyens. Le fait d'être passé au niveau d'alerte 3 augmente le coût opérationnel des forces de l'ordre. Ceci demande donc de la coordination et des budgets plus conséquents.

\section{Les rôles missions etc.}
L'année passée Mr Goergen a eu la tâche de réorganiser le département RH. En effet celui-ci n'était pas un département en tant que tel puisqu'il s'agissait d'une des compétences de la direction générale, au même titre que la logistique ou l'ICT. Il a fallu diviser ces différents départements en entités individuelles.\\

En plus de cela, une loi a été votée incitant [les services publiques \textbf{( vérifier si les services publics aussi)}] et donc la police à faire un pas vers la décentralisation. Ils ont donc choisi de créer 13 entités, sorte de sous-départements responsables d'une zone géographique. La décentralisation d'un compétence est faite si elle apporte une certaine valeur ajoutée. Une partie de celles-ci restent donc centralisées.\\

Le département RH a donc été divisé en 4 entités:
\begin{itemize}
	\item Pour le recrutement et la sélection
	\item Pour la formation avec l'Académie Nationale de Police
	\item Pour la gestion des carrières
	\item Un service psycho-social et médical
\end{itemize}

Alain Georgen a donc mené cette restructuration interne de son département. Il faut maintenant modifier la structure des autres départements de la police.\\

Il a joué plusieurs rôles dans le cadre de ces restructurations. Il s'est définit comme un architecte qui a dut s'assurer du bon fonctionnement des changements, ou encore comme un équilibriste en s'assurant aussi que les changements s'effectuent bien pour les personnes concernées.

%Il a donc géré toute cette partie de modification de la structure, mais dans un autre temps ils vont devoir modifier la structure des autres départements. 
%Avec tous ces changements, ils doivent faire en sorte que tout fonctionne correctement. Il faut coordonner les gens. 
%Il définit donc son rôle comme = un architecte, un équilibriste, toucher à la structure, operationnaliser, et faire en sorte que les choses fonctionnent. (on pourrait le citer). 


%-->Rentrer plus en détail sur ce qui le caractérise comme 
%Dans ce processus de restructuration, Alain Georgen a identifié deux rôles qu'il doit jouer.\\
% Il doit tout d'abord être \textbf{l'architecte de la restructuration:} il doit imaginer, construire et faire en sorte que tout fonctionne. Le succès de la restructuration commencera par le travail que son équipe aura réalisé sur les emplois à délocaliser ou non. \\
% En parallèle à ce rôle d'architecte, il doit jouer aussi celui \textbf{d'équilibriste}. Les pressions de l'état sont assez fortes, lui demandant d'effectuer des coupes budgétaire importantes. De l'autre il doit ne doit pas oublier les personnes qui sont derrière la Police. Il doit être là pour accompagner les personnes lors de ces bouleversements. Des changements de cette ampleur apportent beaucoup d'incertitudes et de crainte pour les employés. La police est composés de personnes avec de grand nombre de fonctions et formations différentes. Il faut donc arriver à faire en sorte que pour chacun des postes, cette restructuration se fera dans les meilleures conditions possibles.\\
% Pour finir, le troisième rôle qu'il doit endosser est celui \textbf{d'agent de changement}. La restructuration commence par le département des ressources humaines. Il doit en quelques sortes montrer l'exemple sur la voie à prendre. 
 %Je ne vois pas vraiment d'ou vient ce paragraphe. 
 
 %--> J'ai peut-être extrapoler ici. Il a mentionné dans l'interview "C’est aussi accompagner le changement parce qu’on est dans une organisation qui est finalement un peu plus fermé que d’autres. On a un turn-over qui est très très bas par rapport à d’autres entreprises." Mais je suis partis en Live a 2h du matin apparement, on peut supprimer donc ^^
\section{Place dans le modèle d'Ulrich}

%\begin{itemize}
%	\item \textbf{Champion des employés:} Assurer que les membres du personnel fassent ce qui est attendu d'eux, y trouvent satisfaction, puissent se développer et rester impliqués 
%	\item \textbf{Agent du changement:} Assurer que la culture de l'organisation évolue de sorte que les comportements des individus soient bien adaptés à l'environnement 
%	\item \textbf{Expert administratif:} Assurer que les règles, processus, systèmes et outils RH fonctionnent de manière efficace et efficiente.
%	\item \textbf{Partenaire stratégique:} Assurer que le déploiement des ressources humaines aide l'organisation à atteindre ses objectifs.
%\end{itemize}

En lien avec le modèle d'Ulrich, on peut commencer par définir Alain Georgen comme un partenaire stratégique. Un partenaire stratégique doit assurer que le de déploiement des réformes mises en place permettent d'atteindre les objectifs recherchés. Cette restructuration a beaucoup de retombées dans différents domaines. A cause de la baisse de budget accordé à la Police, il a fallu redéfinir avec beaucoup d'attention les compétences à décentraliser ou non, prévoir les répercussions financières de celles-ci etc. Ceci a donc nécessité un travail en amont important.\\

Dans la continuité de cette restructuration interne, on peut définir Alain Goergen comme un agent changement. Un agent du changement est quelqu'un qui s'assure que l'évolution de l'entreprise dans laquelle il travaille se passe correctement, que les comportements des individus soient bien adaptés à l'environnement. C'est exactement ce rôle qu'a joue M. Georgen dans le cadre de cette restructuration. 

%Police très éloignée de tous ses employés, sur tout le territoire, le DRH ne voit que ses collègues dans son bureau. C'est donc difficile d'être proche des employés et de pouvoir s'occuper beaucoup d'eux.\\

%Bon expert administratif car tout contrôlé par l'état. Faire très attention à ses dépenses, il faut être le plus efficace possible d'où la restructuration\\
%Je dirai que c'est le rôle du département en général

%--> C'est vrai! Ca peut peut etre faire l'objet d'une réfléxion générale plus tard dans le projet. Chaque DRH s'est impliqué dans une ou des missions spécifiques au cours de la dernière année, néanmoins tous les rôles d'Ulrich sont important pour un bon fonctionnement des RH et ils sont pris en charge par le département ou d'autres personnes

%Je pense que lui ne l'est pas vraiment mais le département oui parce que toute la partie selection/formation est faite pour trouver les talents dont ils ont besoin. 
%aussi un peu de Champion des employés car ils font en sorte de prendre des gens qui ont la bonne mentalité. 
%Par contre je n'arrive pas a retrouver s'il a parlé du fait qu'ils doivent motiver, ....

%--> Même chose que au dessus, est-ce-qu'on parle du département en général dans la partie? C'estun commentaire qu'on peut faire pour les trois DRH je pense et donc peut-être le mettre plus loins dans le travail.

\end{document}


%Comparaison des rôles RH
\section{Comparaison des rôles RH}

Maintenant que nous avons établi les rôles des différents intervenants par rapport à leur discours, nous pouvons engager une comparaison de ceux-ci. \newline

%Comparaison

Pour ce qui est d'Alain Goergen, on voit directement son rôle de gestionnaire du changement à la Police Fédérale. Il s'agit ici d'un rapprochement que l'on peut faire avec Christine Thiran qui doit aussi gérer le changement qui s'opère actuellement chez Mestdagh. Néanmoins, ces deux types de changements sont différents à la Police et chez Mestdagh. M. Goergen a du restructurer le département RH et par la suite les autres départements, ainsi que gérer les inquiétudes de chacun quant à ce changement. Ce changement se fait cependant en incluant les mêmes personnes qui ont déjà adopté la culture organisationnelle de la police. Par contre, dans le cas des magasins Champion, il s'agit d'une fusion avec certains magasins Carrefour. Mme Thiran doit donc gérer les différences d'opinion et les tensions qui peuvent déjà exister entre les deux groupes.\newline \todo[inline]{Parler du fait que Thiran se voit elle-même comme plus proche des gens?=> Motivatrice, bonne humeur etc}

La comparaison du rôle de M. Nolf et M. Goergen est plus simple car ils partagent peu de points communs. En effet, la Police et IBA sont des entreprises bien différentes, qui ont des défis aujourd'hui qui sont bien distincts. IBA se porte bien économiquement ce qui induit qu'elle doit effectuer moins de changements mais plus s'orienter vers une stratégie de croissance alors que la police subit des réductions budgétaires et doit améliorer sa rentabilité avec moins d'effectifs. Ils ne font pas face aux mêmes défis ni au même type de personnel. La même analyse peut se faire entre M. Nolf et Mme Thiran, puisqu'elle travaille actuellement sur un renouvellement de culture organisationnelle, tandis que M. Nolf continue de bâtir une culture organisationnelle déjà existante dans laquelle les nouveaux membres doivent s'intégrer. La première doit recycler deux cultures pour en bâtir une nouvelle, là où le second peut s'appuyer sur une culture existante à améliorer et innover.


%Sorte de conclusion de la comparaison
Suite à cette comparaison, on se rend vite compte que chaque DRH a un rôle bien particulier en fonction de son entreprise, mais aussi et surtout en fonction de la conjoncture de l'entreprise. Nous pouvons encore une fois rappeler que ces DRH ont eu d'autres rôles au cours de leur carrière. 
\todo[inline]{Pas compris la dernière phrase}


%Rôles des DRH dans le modèle d'Ulrich
%Comparaison avec modele ulrich
\section{Rôles des DRH dans le modèle d'Ulrich}

Nous avons à présent une vision des différents rôles que les intervenants endossent. Comme nous le savons, le modèle d'Ulrich donne quatre rôles que peut avoir un responsable de ressources humaines. Il peut donc être intéressant de replacer dans ce modèle, les différents DRH que nous avons rencontré. \\

Commençons par Alain Goergen. Il nous a décrit sa mission comme étant la restructuration complète du département des ressources humaines, et par la suite d'autres départements de la Police Fédérale. Mais il doit aussi aider les employés à s'habituer à ce changement. Au vue de ce rôle qu'il doit jouer, on voit très clairement que cela correspond à l'\og agent de changement \fg{} dans le modèle d'Ulrich. Pour rappel, ce dernier a pour but de faciliter le changement, donc faire en sorte que tout le monde accepte et s'adapte au changement.\\

%En lien avec le modèle d'Ulrich, on peut commencer par définir Alain Georgen comme un partenaire stratégique. Un partenaire stratégique doit assurer que le de déploiement des réformes mises en place permettent d'atteindre les objectifs recherchés. Cette restructuration a beaucoup de retombées dans différents domaines. A cause de la baisse de budget accordé à la Police, il a fallu redéfinir avec beaucoup d'attention les compétences à décentraliser ou non, prévoir les répercussions financières de celles-ci etc. Ceci a donc nécessité un travail en amont important.\\

%Dans la continuité de cette restructuration interne, on peut définir Alain Goergen comme un agent changement. Un agent du changement est quelqu'un qui s'assure que l'évolution de l'entreprise dans laquelle il travaille se passe correctement, que les comportements des individus soient bien adaptés à l'environnement. C'est exactement ce rôle qu'a joue M. Georgen dans le cadre de cette restructuration. 

Nous pouvons aussi le voir comme \og partenaire stratégique \fg{}. En effet, la décentralisation qu'il doit opérer, a un impacte stratégique en plus d'un impact sur les employés comme présenté ci-dessus. En reprenant les mots de Mr Goergen \og Le modus était tout ce qui peut être déconcentré doit être déconcentré mais si il n’y a pas de plus value à déconcentrer il faut le garder au niveau central \fg{} \footnote{TO DO : url video + minute}. Il parle ici de \og plus-value \fg{}, il y a donc fait une analyse stratégique sur les compétences à décentraliser. Dans ce cas, il ne lui a pas été demandé d'appliquer un plan pensé en amont, mais de le créer par lui même. Ses décisions ont donc été prises en fonction de l'objectif d'efficacité à atteindre. Il est aussi important de se rappeler que la Police fait fasse à des réductions de budget et que ses décisions doivent les prendre en compte. Cela fait donc de lui un \og partenaire stratégique \fg{} au sens d'Ulrich.\\   

\todo[inline]{This is a shaft}

Pour ce qui est de Frédéric Nolf, GRH d'IBA, on peut le classer lui aussi en tant que \textbf{agent de changement} dans le modèle d'Ulrich. Le rôle qu'il a joué durant la constitution d'une dream team pour le nouveau CEO permet de le classer assez facilement dans cette catégorie. Mais ce qui conforte ceci est la façon dont Nolf profile les DRH de nos jours:\newline

\og{} \textit{Je pense que les DRH modèles [\ldots] se profilent tous comme 
des architectes et des transformateurs d’entreprise \ldots} \fg{} \newline

Mais Nolf correspond aussi à un \textbf{champion des employés} dans le modèle d'Ulrich dans la façon dont il traite le personnel d'IBA. En effet, il prend en compte les avis des gens afin que ces derniers ne se sentent pas lésés et restent motivés. De plus, la création d'activités communes comme le sport ou bien de cellules de réflexion sur le développement durable permettent aux employés de s'épanouir au sein de la société. On a donc une conception très humaine de l'employé. On veut que chaque ingénieur et chercheur se sente bien dans la société. \newline

Enfin, on peut considérer le GRH d'IBA comme un \textbf{partenaire stratégique}. Un des soucis de Frédéric Nolf, comme indiquer plus haut, est de garder les employés concentrés dans leurs tâches. Les chercheurs sont des gens qui font facilement des choses à leur manière. De plus, une autre de ses préoccupations est d'exploiter le potentiel de chaque employé malgré l'internationalisation. Tout ceci est dans l'optique d'augmenter la productivité, l'efficacité de l'entreprise. Produire plus mais avec le même effectif. 



%Conclusion du la comparaison au modèle 


%Limites de la théorie
%Partie 4 A : Prise de recul et discussion des limites de la théorie
\subsection{Limites de la théorie}

Commençons donc par détailler ce que nous pensons être les limites de ce modèle. Les différents points que nous aborderons sont des rôles multiples que peut avoir un DRH, du manque de temporalité et pour finir, du département RH.\

\subsubsection*{Multiplicité des rôles} 

Le premier point qui saute aux yeux, lors du rapprochement entre les rôles de ces DRH et le modèle d'Ulrich, est la possibilité d'avoir plusieurs rôles. \newline

En se basant sur son discours, on observe que Christine Thiran correspond à l'\og{}agent de changement\fg{}, de par sa mission d'accompagnement de la fusion du groupe Mestdagh avec des magasins Carrefour, mais aussi comme \og{}championne des employés\fg{}. \newline

M. Nolf dispose lui aussi de plus d'un rôle. Le premier qu'on peut retrouver est d'être le \og{}champion des employés\fg{}. Il a pour rôle de répondre aux besoins et aux préoccupations des employés mais également d'assurer leur motivation et leur contribution aux objectifs de l'entreprise. Parallèlement, on peut aussi l'associer au rôle de \og{}partenaire stratégique \fg{} puisqu'il a pour mission de trouver les talents de demain et d'assurer que ceux déjà présents permettent à l'organisation d'atteindre ses objectifs, en collaboration avec les autres dirigeants. Dans une moindre mesure (en tout cas, moins discutée lors de la conférence), il est également \og{}agent de changement\fg{} puisqu'il fait évoluer les comportements au sein de l'entreprise pour qu'ils collent à la culture IBA et également qu'IBA suive les avancées technologiques de manière continue. 

Ces personnes assument donc plusieurs rôles du modèle théorique. \\

Il s'agit donc ici d'une première limite du modèle d'Ulrich. Chaque rôle présent dans ce modèle correspond à une période de l'histoire avec son contexte économique, une manière de concevoir l'organisation, une conception de la main d'\oe{}uvre. Chacun d'eux exprime ce que fait le responsable des ressources humaines dans son entreprise. Ce modèle donne donc un rôle à chaque DRH en fonction de sa mission, à un cmoment donné, mais ne permet pas qu'il puisse en avoir plusieurs. Dave Ulrich lui-même considère aujourd'hui (lors de ses conférences et dans ses écrits) qu'un leader est naturellement bon dans 2 des 4 cadrans, mais le travail aujourd'hui est de le faire devenir bon dans les 4 cadrans. (Dave Ulrich, 2013, \cite{leadership2013})\newline 

La première limite de ce modèle est donc le manque d'une pluralité des fonctions pour un même DRH. En réalité, on sait que nous avons hérité de certaines pratiques de chacun des rôles du modèle d'Ulrich (Nathalie Delobbe, 2015,\cite{slidegrh}). Il est donc tout à fait possible qu'aujourd'hui un responsable des ressources humaines assume plusieurs fonctions en même temps dans son métier.
\footnote{Comme cela a été fait dans une étude qui assignait un pourcentage de chaque rôle en fonction du temps sur plusieurs périodes par Lawler, Boudreau \& Mohrman (Lawler, 2006, \cite{Lawler2006})}



\subsubsection*{Historicité des rôles}

Le deuxième point que nous souhaitons aborder dans les limites de ce modèle d'Ulrich, est le manque de temporalité, ou encore, de dynamisme dans le modèle. Certes, comme dit précédemment, ce modèle a été construit en fonction de l'évolution du rôle de la fonction RH dans le temps. Mais, par contre, il n'exprime pas la possibilité d'une évolution entre les rôles, pour un même DRH. Ce point est proche du précédent. En effet, il s'agit aussi du fait qu'un DRH puisse avoir plusieurs rôles. Mais cette fois nous attirons l'attention sur le point de vue dynamique du rôle. \\

Les trois intervenants nous ont présenté leur parcours en tant que responsable des ressources humaines. Ces parcours, bien que fort différents, ont un point commun : tous ont eu plusieurs rôles au fil du temps. Prenons par exemple Alain Goergen, qui dans la même entreprise est passé par les ressources humaines dans le recrutement, la formation et aujourd'hui dans la réorganisation complète de son département. On voit donc une évolution de sa mission au cours du temps. Selon le modèle d'Ulrich, son rôle a évolué dans le temps, passant par exemple du rôle d'\og{}expert administratif\fg{} à celui d'\og{}agent de changement \fg{}. \newline 

Cette évolution du rôle des DRH dans le temps, est un autre manquement au modèle théorique. C'est ce que nous avons appelé ici plus haut le manque de temporalité, ou encore, de dynamisme. Il est donc important de garder à l'esprit que dans une société, une personne a bien souvent un rôle qui évolue au cours du temps, en fonction du contexte économique, social ou opérationnel de son entreprise. \


\subsubsection*{Du rôle du département RH}

Le troisième et dernier point que nous aborderons ici, comme limite de ce modèle théorique, est le département RH en lui même. Depuis quelques temps déjà, les entreprises se sont rendues compte de l'importance de leurs employés, et de leur bien-être. Les départements sont apparus et ils prennent une place de plus en plus importante dans les sociétés. Les intervenants nous l'ont dit, ils sont de plus en plus proches de la direction et leur département est de plus en plus grand, d'un côté parce que la taille des entreprises croît mais aussi au vu de l'importance qui y est accordée. \newline
\todo[inline]{Une source serait bien pour montrer que les entreprises se rendent compte de l'importance de leurs employés}

L'évolution de la taille de ce département, pour quelle raison que ce soit, permet de couvrir plus de missions. Le département peut s'occuper de plus de choses au quotidien. Nous parlons donc encore une fois du fait d'assumer plusieurs rôles à la fois, au sens d'Ulrich. Mais il ne s'agit plus ici d'une seule personne, il s'agit du rôle du département entier, l'ensemble des personnes qui le composent. Le plus bel exemple est sûrement celui de la Police Fédérale. Alain Goergen fait en effet partie d'un département de plus de 500 personnes. Comme il nous l'a présenté, ce département est découpé en entités qui ont chacune un rôle (recrutement, formation, psycho et socio-médical, ...). Ils recouvrent en fait tous les rôles présents dans le modèle d'Ulrich, même si chaque personne n'a peut-être qu'un seul rôle. \newline

Ce troisième point découle donc du premier puisqu'il s'agit aussi de la multitude de rôles de la fonction RH, cette fois pas pour une seule personne, mais sur le département entier. \newline

%Conclusion
Enfin, si nous devions recréer un nouveau modèle pour les rôles du DRH, il faudrait surement prendre en compte les trois limites citées ci-dessus. Même si elles ne se basent que sur l'analyse de trois DRH, on peut supposer qu'il en va de même pour beaucoup d'autres. Le modèle d'Ulrich n'est donc plus exactement à jour pour analyser les ressources humaines d'aujourd'hui, mais il est évident qu'il ne faut pas l'écarter pour autant car il permet facilement d'avoir une visibilité du positionnement et de l'action des RH. Il nous permet également de voir comment a évolué cette fonction dans les entreprises, au cours de l'histoire, mais aussi d'où viennent les pratiques que les responsables de ressources humaines actuels appliquent tous les jours. \\


\bibliographystyle{plain}
\bibliography{biblio}

\end{document}
